% =====================================================================================
%  Chapter 6 — Two-Robot Experimental Results
%  Thesis: Comparison of Control Strategies for Interaction-Force Aware Multirotor Teams
%
%  DRAFT SKELETON, 27 September 2026.
%  Sources: docs/30_Results_Chapters_Skeleton.md (§6.1--6.6), docs/05, docs/10, docs/25,
%           Chapter~\ref{ch:setup} (platform and protocol).
%
%  HONESTY NOTE. This file is structure and framing only. No C.4 comparison numbers, figures,
%  or rankings appear here. Every slot for measured outcomes is marked with \todo so the compiled
%  PDF cannot be read as reporting results that do not yet exist.
%
%  Citation baseline: 2026-08-27 (see docs/thesis/CITATION_BASELINE.md).
% =====================================================================================

\chapter{Two-Robot Experimental Results}
\label{ch:results2}

This chapter reports the systematic comparison on the two-vehicle team once the C.4 campaign is
complete. Chapter~\ref{ch:setup} fixed the platform, scenarios and protocol; Chapter~\ref{ch:control}
defined Strategies~0--3 compared against the baseline. The sections below follow the analysis
pipeline from aggregate tracking metrics through control effort, the predictive residual path, and a
summary of when each strategy is preferable. Research questions RQ1, RQ2 and RQ4
(Section~\ref{sec:intro-rq}) are answered primarily here; RQ3 is deferred to
Chapter~\ref{ch:results3}.

% -------------------------------------------------------------------------------------
\section{Experimental campaign overview}
\label{sec:res2-overview}

The comparison flights use the brushless pair \textbf{cf5} (study vehicle, bottom role in vertical
stacks) and \textbf{cf\_second} (top), with external pose from a four-camera OptiTrack system and
onboard logging at \SI{500}{\hertz} on micro-SD as the source of record (Section~\ref{sec:setup-platform}).
Radio telemetry is for monitoring only; thesis metrics are computed from merged uSD logs after role
alignment, using the same pipeline documented in Chapter~\ref{ch:setup}.

Scenarios are drawn from the frozen formation library (Chapter~\ref{ch:setup}, Section~\ref{sec:setup-scenarios}):
the C.1/C.4 matrix prioritises vertical separation and downwash-dominated geometries (e.g.\ A1,
A3, A4, A7) together with null coplanar controls (C1--C3) that bound the measurement floor. Each
(scenario, strategy) cell is intended to be repeated $n$ times under interleaved runtime controller
selection so session effects do not align with strategy choice (Table~\ref{tab:setup-threats}).

The compared set follows Chapter~\ref{ch:control}: Strategy~0 (uncompensated geometric baseline),
Strategy~1 (pure INDI), Strategy~2 (geometric position loop with onboard Neural-Swarm2-style
residual prediction), and Strategy~3 (feedback linearisation with the same predictor), each flown
with gains frozen after the hardware gate. Indicative single flights and partial C.1 collection
exist; the tables and figures in the following sections are reserved for the completed C.4 repeats.

% -------------------------------------------------------------------------------------
\section{Tracking and attitude}
\label{sec:res2-tracking}

This section answers the tracking half of RQ1: whether compensating strategies reduce position and
separation error relative to Strategy~0 under identical trajectories, and how reactive versus
predictive families rank. Attitude excursions matter for safety and for interpreting residual
magnitude; they are reported alongside position error rather than as a separate optimisation target.

\todo[color=red!25]{PLACEHOLDER: Table~\ref{tab:res2-tracking} populated once C.4 flights complete.
Columns: scenario, strategy, RMSE $x$/ $y$/ $z$/ $\|\mathbf{e}\|$, peak position error, peak
separation error, attitude peak (roll/pitch), $n$ repeats. Rows: all (scenario, strategy) cells in
the comparison matrix. Statistics: mean $\pm$ std across repeats.}

\begin{table}[t]
  \centering
  \caption{Tracking error summary (two-robot C.4 campaign).}
  \label{tab:res2-tracking}
  \small
  \begin{tabular}{@{}llrrrrr@{}}
    \toprule
    Scenario & Strategy & RMSE $x$ & RMSE $y$ & RMSE $z$ & Peak $\|\mathbf{e}\|$ & $n$ \\
    \midrule
    \multicolumn{7}{@{}l@{}}{\textit{--- populated from phase metrics / \texttt{plot\_comparison.py} ---}} \\
    \makecell{Placeholder} & S0 & --- & --- & --- & --- & --- \\
    \bottomrule
  \end{tabular}
\end{table}

\todo[color=red!25]{PLACEHOLDER: Figure (grouped bars) --- controllers $\times$ metrics with SEM
error bars; source \texttt{plot\_comparison.py} on phase CSV. One panel per metric family (position
RMSE, peak error).}

\todo[color=red!25]{PLACEHOLDER: Figure (time series) --- commanded vs achieved position or
separation for one representative scenario (e.g.\ A8 or A4), overlay Strategies~0--3 on the same
axes; source \texttt{plot\_flight.py} or dashboard export.}

A short subsection will contrast geometric tracking (Strategy~0), our INDI implementation
(Strategy~1), and any reference INDI flights included for context, without treating reference
controllers as part of the primary hypothesis test unless explicitly added to the C.4 matrix.

% -------------------------------------------------------------------------------------
\section{Control effort and safety}
\label{sec:res2-effort}

Better tracking is not automatically preferable if it requires sustained saturation or leaves less
margin on separation. This section supports RQ4 by pairing performance with effort proxies and
closest-approach statistics.

\todo[color=red!25]{PLACEHOLDER: Table~\ref{tab:res2-effort} --- thrust/torque proxies, motor
saturation fraction (\texttt{effort\_proxy}, \texttt{motor\_*} from phase metrics), minimum
horizontal separation (\texttt{min\_sep\_m}), minimum vertical offset (\texttt{min\_dz\_m}), mean
$\pm$ std, $n$.}

\begin{table}[t]
  \centering
  \caption{Control effort and safety margins (two-robot C.4 campaign).}
  \label{tab:res2-effort}
  \small
  \begin{tabular}{@{}llrrrr@{}}
    \toprule
    Scenario & Strategy & Effort proxy & Sat.\ fraction & Min sep.\ (\si{\metre}) & $n$ \\
    \midrule
    \makecell{Placeholder} & S0 & --- & --- & --- & --- \\
    \bottomrule
  \end{tabular}
\end{table}

The prose here will state explicitly when a strategy trades higher effort for lower tracking error,
and when reduced separation margin disqualifies a configuration despite good RMSE.

% -------------------------------------------------------------------------------------
\section{Predictive strategy (Neural-Swarm2)}
\label{sec:res2-nn}

Strategy~2 is evaluated in two layers: open-loop agreement between predicted and measured
$a_{\mathrm{res}}$ on the C.1 training distribution, and closed-loop C.4 flights with prediction
enabled versus Strategy~0 on matched scenarios. RQ1 and RQ2 both depend on the second; the first
documents model quality independent of loop gain.

\todo[color=red!25]{PLACEHOLDER: Figure --- predicted vs measured $a_{\mathrm{res}}$ (scatter,
$R^2$ or correlation); second panel error vs relative geometry (distance, $\Delta z$, gated flag).
Source: merged uSD + \texttt{firmware\_forward} replay where applicable.}

\todo[color=red!25]{PLACEHOLDER: Table --- Strategy~2 vs S0 on matched scenarios (same columns as
Table~\ref{tab:res2-tracking}) for in-flight enable; note training scenarios and weight checkpoint
(\texttt{full\_bank\_*.npz} provenance).}

% -------------------------------------------------------------------------------------
\section{Hybrid (NA-INDI), if in scope}
\label{sec:res2-hybrid}

The primary compared set in Chapter~\ref{ch:control} is Strategies~0--3. A NA-INDI / LINDI hybrid
(Strategy~4 in the wider project numbering) is \emph{not} in that set unless explicitly restored to
the C.4 matrix. This section remains as a structured slot: if hybrid flights are flown, they use
the same metric tables as Sections~\ref{sec:res2-tracking}--\ref{sec:res2-effort} for fair
comparison.

\todo[color=red!25]{PLACEHOLDER (conditional): Configuration block --- $c$ parameter, weight source
(retrain vs reference port). Table shell identical to §\ref{sec:res2-tracking}. Omit section from
final PDF if hybrid remains out of scope.}

% -------------------------------------------------------------------------------------
\section{Summary --- when which strategy}
\label{sec:res2-summary}

This section closes the two-robot story for RQ1 and RQ2: a scenario-class breakdown (vertical stack,
lateral crossing, near-ground, fast transient) with both qualitative ranking and the quantitative
backing from Tables~\ref{tab:res2-tracking}--\ref{tab:res2-effort}.

\todo[color=red!25]{PLACEHOLDER: Summary table --- rows: strategies; columns: scenario classes;
cells: best tracking / effort trade-off / data \& compute cost (RQ4 pointers). No invented
rankings until C.4 data exist.}

\begin{table}[t]
  \centering
  \caption{Strategy suitability by scenario class (to be filled after C.4).}
  \label{tab:res2-when-which}
  \small
  \begin{tabular}{@{}lcccc@{}}
    \toprule
    Strategy & Vertical stack & Lateral traverse & Near-ground & Fast transient \\
    \midrule
    S0 (baseline) & --- & --- & --- & --- \\
    S1 (INDI) & --- & --- & --- & --- \\
    S2 (pred.\ + geo) & --- & --- & --- & --- \\
    S3 (pred.\ + FBL) & --- & --- & --- & --- \\
    \bottomrule
  \end{tabular}
\end{table}

Train/test splits by scenario class (Section~\ref{sec:setup-validity}) feed RQ2 directly: a crossover
in favour of reactive or predictive compensation should appear here if it exists in the data.
